\documentclass[12pt]{article}

\usepackage[utf8]{inputenc}
\usepackage[T1]{fontenc}
\usepackage[spanish]{babel}
\usepackage{geometry}
\usepackage{setspace}

\geometry{a4paper, margin=1in}
\setstretch{1.3}
\setlength{\parskip}{0.6em}
\setlength{\parindent}{0pt}

\begin{document}

\begin{center}
{\Large\bfseries Open-Vocabulary Semantic Change Detection in Bitemporal Hyperspectral Images}\\[0.5cm]
{\large Detección Semántica de Cambios con Vocabulario Abierto en Imágenes Hiperespectrales Bitemporales}
\end{center}
\vspace{0.6cm}

\section*{Objetivo General}

Desarrollar un sistema computacional para la detección y caracterización semántica de cambios con vocabulario abierto en imágenes hiperespectrales bitemporales, mediante la integración de modelos fundacionales, técnicas de segmentación guiada por indicaciones y modelos de visión-lenguaje y generación de descripciones textuales.

\section*{Objetivos Específicos}

\begin{enumerate}
    \item Extraer representaciones profundas de las imágenes hiperespectrales bitemporales mediante un modelo fundacional y obtener, a partir de dichas representaciones, un mapa inicial de cambios que permita localizar las regiones con diferencias significativas entre las dos adquisiciones.
    \item Integrar un modelo de segmentación guiado por indicaciones espaciales derivadas del mapa inicial de cambios, para delimitar a nivel de píxel las regiones de cambio.
    \item Incorporar un modelo de visión-lenguaje para identificar y caracterizar las regiones de cambio mediante etiquetas semánticas de vocabulario abierto, a partir de la similitud entre las representaciones de las regiones y diferentes descripciones textuales de posibles tipos de cambio.
    \item Generar descripciones en lenguaje natural de las regiones de cambio a partir de sus características espectrales, con el propósito de complementar las etiquetas de vocabulario abierto con información semántica más detallada.
    \item Evaluar cuantitativamente el desempeño del sistema sobre un conjunto de datos hiperespectral bitemporal, utilizando métricas de detección de cambios y analizar la capacidad del sistema para caracterizar semánticamente las regiones detectadas.
\end{enumerate}

\section*{Justificación}

En los últimos años, los modelos fundacionales de visión por computador han transformado las tareas de segmentación y de comprensión semántica de imágenes. Los modelos de segmentación guiados por indicaciones (prompts) como Segment Anything, conocido como SAM, permiten obtener máscaras precisas de objetos sin haber sido entrenados para categorías específicas [1], mientras que los modelos de visión y lenguaje como CLIP habilitan la clasificación de vocabulario abierto al alinear imágenes y texto en un mismo espacio semántico [2]. Sobre esta base, la detección semántica de cambios con vocabulario abierto, denominada en inglés Open-Vocabulary Semantic Change Detection, ha comenzado a consolidarse en teledetección: métodos recientes como Semantic-CD identifican y describen los cambios entre imágenes de una misma escena sin limitarse a un conjunto cerrado de clases definidas de antemano [3]. No obstante, todos estos desarrollos operan sobre imágenes ópticas RGB o multiespectrales; para imágenes hiperespectrales, que registran cientos de bandas espectrales por píxel, su extensión constituye un problema poco explorado..

Llevar la detección semántica de cambios con vocabulario abierto al dominio hiperespectral abre un abanico de aplicaciones de alto impacto. La firma espectral de cada píxel actúa como una huella dactilar del material, de modo que un sistema capaz de detectar y describir cambios sobre este tipo de imágenes podría distinguir transformaciones imperceptibles para una cámara convencional. En el ámbito civil, esto resulta valioso para el monitoreo de la deforestación, la minería ilegal, la contaminación de cuerpos de agua, el crecimiento urbano, la salud de los cultivos y la respuesta ante desastres naturales, escenarios en los que la naturaleza del cambio importa tanto como su ubicación. En el ámbito militar y de seguridad, la riqueza espectral permite revelar materiales y actividades que el camuflaje visual oculta, como movimientos de vehículos, alteraciones del terreno o presencia de materiales artificiales, información difícil de obtener con sensores RGB tradicionales.

Para abordar esta brecha, el presente proyecto buscará integrar un conjunto de modelos fundacionales ya existentes en una única cadena de procesamiento: extracción de características espectrales y espaciales de las imágenes, localización de cambios, segmentación de regiones y asignación de etiqueta de vocabulario abierto con la posibilidad de generar descripciones en lenguaje natural mediante modelos de generación de descripciones espectrales (spectral captioning).

Con base en lo anterior, este proyecto se orienta por la siguiente pregunta de investigación: ¿Cómo realizar detección semántica de cambios con vocabulario abierto en imágenes hiperespectrales bitemporales mediante la alineación entre información espectral y textual? Responder a esta pregunta implica diseñar una solución que, sin entrenamiento supervisado específico para cada tipo de cambio, sea capaz de localizar, segmentar y describir con vocabulario abierto las diferencias entre dos adquisiciones hiperespectrales de una misma escena.

El sistema propuesto se estructura como una cadena de procesamiento secuencial. Inicialmente se carga el par de imágenes hiperespectrales de la misma escena y se alinean a un espacio espectral común [4]. Se extraen representaciones profundas de cada imagen. Luego se calcula un mapa de cambio comparando ambas representaciones y se binariza para aislar las regiones modificadas. Posteriormente, esas regiones se convierten en indicaciones espaciales, que producen máscaras precisas a nivel de píxel. Finalmente, se etiqueta cada máscara contra un vocabulario textual libre y, de forma opcional, se genera una descripción en lenguaje natural de cada región con un modelo de generación de descripciones espectrales [5],[6]. Como resultado, el sistema generaría para cada par de imágenes un conjunto de regiones de cambio, cada una acompañada de su máscara, su etiqueta semántica y, cuando corresponda, una descripción textual del cambio detectado.

\begin{thebibliography}{99}

\bibitem{ref1} A. Kirillov, E. Mintun, N. Ravi, H. Mao, C. Rolland, L. Gustafson, T. Xiao, S. Whitehead, A. C. Berg, W.-Y. Lo, P. Dollár, and R. Girshick, ``Segment anything,'' in \textit{Proceedings of the IEEE/CVF International Conference on Computer Vision (ICCV)}, pp. 4015--4026, 2023. arXiv:2304.02643.

\bibitem{ref2} A. Radford, J. W. Kim, C. Hallacy, A. Ramesh, G. Goh, S. Agarwal, G. Sastry, A. Askell, P. Mishkin, J. Clark, G. Krueger, and I. Sutskever, ``Learning transferable visual models from natural language supervision,'' in \textit{Proceedings of the 38th International Conference on Machine Learning (ICML)}, vol. 139 of PMLR, pp. 8748--8763, 2021.

\bibitem{ref3} Y. Zhu, L. Li, K. Chen, C. Liu, F. Zhou, and Z. Shi, ``Semantic-CD: Remote sensing image semantic change detection towards open-vocabulary setting,'' \textit{arXiv preprint arXiv:2501.06808}, 2025.

\bibitem{ref4} D. Wang, M. Hu, Y. Jin, Y. Miao, J. Yang, Y. Xu, X. Qin, J. Ma, L. Sun, C. Li, C. Fu, H. Chen, C. Han, N. Yokoya, J. Zhang, M. Xu, L. Liu, L. Zhang, C. Wu, B. Du, D. Tao, and L. Zhang, ``HyperSIGMA: Hyperspectral intelligence comprehension foundation model,'' \textit{IEEE Transactions on Pattern Analysis and Machine Intelligence}, 2025. arXiv:2406.11519.

\bibitem{ref5} J. Li, D. Li, C. Xiong, and S. C. H. Hoi, ``BLIP: Bootstrapping language-image pre-training for unified vision-language understanding and generation,'' in \textit{Proceedings of the 39th International Conference on Machine Learning (ICML)}, vol. 162 of PMLR, pp. 12888--12900, 2022. arXiv:2201.12086.

\bibitem{ref6} A. Das, T. Rachamalla, P. Singh, K. Biswas, V. K. Verma, S. Garcia, A. Plaza, and S. K. Roy, ``HyperCap: Hyperspectral land cover captioning dataset for vision language models,'' \textit{IEEE Geoscience and Remote Sensing Magazine}, 2026. arXiv:2505.12217.

\end{thebibliography}

\end{document}
